\documentclass[journal,10pt]{IEEEtran}
\usepackage{amsmath,amssymb}
\usepackage{booktabs}
\usepackage{array}
\usepackage{graphicx}
%\usepackage{balance}   % 暂时停用：末页双栏平衡在内容变短后产生 1.7pt 竖直溢出
\usepackage[dvipsnames]{xcolor}
\usepackage[colorlinks=true,linkcolor=NavyBlue,citecolor=PineGreen,urlcolor=NavyBlue]{hyperref}
\usepackage{xspace}
%% 由 scripts/make_tables.py 生成，勿手改；取值来自 results/c5_matrix.json 与 results/r37e_full_seeds.json。
%% 正文与表说明里的这些数字只允许写成这些宏。
\newcommand{\cFiveAllSparseYellowA}{547}
\newcommand{\cFiveEnergyDeadA}{0}
\newcommand{\cFiveEnergyDeadB}{0}
\newcommand{\cFiveEnergyLostA}{290}
\newcommand{\cFiveEnergyLostB}{826}
\newcommand{\cFiveEnergyYellowA}{559}
\newcommand{\cFiveEnergyYellowB}{1194}
\newcommand{\cFiveNightfloorDeadA}{20}
\newcommand{\cFiveNightfloorDeadB}{31}
\newcommand{\cFivePeak}{0.012}
\newcommand{\cFiveSeeds}{3}
\newcommand{\cFiveTtlEightDeadA}{0}
\newcommand{\cFiveTtlEightDeadB}{0}
\newcommand{\cFiveTtlEightRevertA}{10}
\newcommand{\cFiveTtlEightRevertB}{9}
\newcommand{\cFiveTtlEightYellowA}{849}
\newcommand{\cFiveTtlEightYellowB}{2020}
\newcommand{\cFiveTtlFourLostB}{104}
\newcommand{\cFiveTtlFourRevertA}{6}
\newcommand{\cFiveTtlFourRevertB}{5}
\newcommand{\cFiveTtlFourYellowA}{849}
\newcommand{\cFiveTtlFourYellowB}{1916}
\newcommand{\cFiveValidUntilYellowA}{849}
\newcommand{\cFiveValidUntilYellowB}{2020}
\newcommand{\cFiveYellowNA}{2016}
\newcommand{\cFiveYellowNB}{3528}
\newcommand{\cThreeDeathsFifo}{12}
\newcommand{\cThreeDeathsPurge}{0}
\newcommand{\cThreeExpiredFifo}{2538}
\newcommand{\cThreeExpiredPurge}{120}
\newcommand{\cThreeLatestDelta}{5.55}
\newcommand{\cThreePairedHi}{4.60}
\newcommand{\cThreePairedLo}{3.48}
\newcommand{\cThreePairedPoints}{4.04}
\newcommand{\cThreeRatio}{2.45}
\newcommand{\cThreeSeedZeroFifo}{202}
\newcommand{\cThreeSeedZeroPurge}{420}
\newcommand{\cThreeSvcFifo}{0.373}
\newcommand{\cThreeSvcPurge}{0.413}
\newcommand{\cThreeTenSeedFifo}{1691}
\newcommand{\cThreeTenSeedPurge}{4145}
   % 由 scripts/make_tables.py 生成，勿手改
\raggedbottom   % 末页不满时按自然高度结束，避免 1pt 级的竖直溢出
\newcommand{\stime}{$S_{\mathrm{time}}$\xspace}
\newcommand{\scap}{$S_{\mathrm{cap}}$\xspace}
\newcommand{\senergy}{$S_{\mathrm{energy}}$\xspace}
\newcommand{\saccess}{$S_{\mathrm{access}}$\xspace}
\newcommand{\unk}{$U$\xspace}
\newcommand{\trel}{t_{\mathrm{rel}}}

\title{\textbf{Local Communication Control for Pre-Disaster Monitoring\\
with Intermittent Backhaul}}

\author{Pre-disaster monitoring \& emergency-communication working draft\\
\texttt{adaptive-communication} repository --- working draft}

\date{}

\begin{document}
\maketitle

\begin{abstract}
Intermittent backhaul disconnects a monitoring controller without stopping the communication
state it has installed. Unacknowledged records continue consuming access opportunities, and
applied sampling profiles continue drawing energy after corrective commands become unreachable.
We study this failure in pre-disaster mountain monitoring with battery/solar LoRaWAN Class~A
nodes, a cellular primary and an uplink-only short-message backup. We formulate an execution
contract that separates record lifetime, authorised configuration fallback and evidence of
end-to-end completion, and examine where each decision can be enforced. The central observation
is that a valid local action can require less information than complete failure diagnosis.
In the tested cache model, placing standard deadline expiry at the source improves on-time
service by \cThreePairedPoints\ percentage points over ten paired seeds and increases outage
on-time delivery by \cThreeRatio$\times$; gateway-only suppression can instead back-pressure
access. Node-local protection also prevents losses that a central downgrade cannot avert in
the tested episodes. Ordinary TTL and battery rules cover the evaluated configuration-stopping
problem, including a same-information exact reference, so a new lease algorithm is not required
for these results. Partial gateway attribution and live Agent traces expose a separate interface
failure: access receipts being interpreted as backhaul or completion evidence. These findings
support local enforcement and explicit evidence boundaries as communication-system design
principles. The Agent study is formative; an independent benefit of the complete execution
interface remains to be evaluated against a same-component ordinary stack.
\end{abstract}

\begin{IEEEkeywords}
Pre-disaster monitoring, intermittent backhaul, LoRaWAN Class~A, communication state,
local enforcement, deadline expiry, partial observability, agent execution.
\end{IEEEkeywords}

\section{Introduction}
A mountain monitoring node receives an authorised dense sampling profile while its backhaul is
available. The cellular path then fails, but the node continues sampling and transmitting. A
later downgrade cannot reach its Class~A receive window. Meanwhile, records lacking central
acknowledgement remain in the source cache and occupy later access batches even after their
business deadlines. The gateway may suppress useless return traffic without releasing the
source copies. Correct decisions at one layer can therefore fail to stop resource consumption
at another.

This paper studies \emph{persistent communication state under disconnected control}. Record
retention and applied configuration have different termination conditions, but share an execution
dependency: waiting for a later central decision assumes a control path that may be unavailable.
The monitoring obligation remains end-to-end---an in-window sample must arrive at the centre
before its deadline---while evidence and authority are distributed across the node, gateway and
centre. An access receipt confirms one stage of that chain, not the entire obligation.

\subsection{Design insight}
\textbf{Local action and complete diagnosis have different information requirements.} A gateway
may be unable to distinguish a sample that was never collected from one not yet heard. The
source nevertheless knows whether a stored record has passed its locally assigned deadline and
can stop retransmitting it. Similarly, an authorised local fallback can act on a clock or
battery condition even when the centre cannot determine the current profile. These actions
need explicit authority and enforceable local predicates; they do not need an invented account
of the remote failure.

We organise this distinction as an execution contract: identify the state an operation changes,
the evidence supporting its outcome, and the location capable of acting before further harm.
Standard expiry and local protection instantiate the contract. The contribution is evaluated
through their placement and communication consequences, rather than through a new name for
these established mechanisms.

\subsection{Ordinary mechanisms and the remaining question}
The comparisons include standard per-record expiry, local clock and battery protection,
fixed-duration fallback, strong backup packing and a same-information exact stopping reference.
Expiry reproduces obligation-keyed deletion on the tested seeds, and ordinary configuration
rules cover the evaluated stopping family. These results delimit the need for a new algorithm.
They do not establish a global optimum for communication scheduling.

A remaining systems question is whether an explicit execution interface improves action selection
and repair beyond a stack already containing those ordinary components. The distinction matters
for Agents: a planner may order more sampling when the missing stage is return delivery, or infer
that a command applied from unrelated access receipts. Our live study observes the latter
interface failure, but its end-to-end comparison does not yet isolate the interface's independent
contribution. We keep that question separate from the deterministic placement results.

\subsection{Contributions}
\begin{enumerate}
\item \textbf{An execution contract for disconnected monitoring}
(\S\ref{sec:model}--\S\ref{sec:mech}). The formulation separates locally decidable record
expiry, authorised configuration fallback and partial end-to-end evidence. An enforcement
condition combines evidence availability, authority and actuation delay.
\item \textbf{Measured cross-segment placement effects}
(\S\ref{sec:eval-record}, \S\ref{sec:eval-place}). Source-local standard expiry improves
service by \cThreePairedPoints\ percentage points (95\% CI
[\cThreePairedLo,\cThreePairedHi]) and outage on-time delivery by \cThreeRatio$\times$.
A complementary configuration experiment shows why a central protection rule can fail when
the same rule can still act locally.
\item \textbf{An empirical boundary for control and evidence}
(\S\ref{sec:eval-lease}--\S\ref{sec:eval-agent}). Same-information ordinary stopping rules
match the declared exact reference; gateway attribution remains necessarily partial; and
formative Agent traces identify access/backhaul evidence confusion. Together these comparisons
separate established local remedies from unproven planner or runtime increments.
\end{enumerate}
The setting is pre-disaster monitoring with externally authorised requirements. The system does
not infer hazard levels, remove missed obligations or obtain additional communication resources
by changing the scoring task.

\section{Scenario and Operating Point}
\label{sec:bg}

Fourteen sensor nodes in two groups monitor displacement and rainfall on a slope. Nodes sample
and upload to a local gateway over LoRaWAN Class~A; the gateway forwards to a central server over
cellular backhaul, with a BeiDou short-message service as a sparse, metered backup carrying
field-to-centre messages only, never downlink configuration. Nodes are battery plus small solar,
the mission is pre-disaster and routine with authorised warning-level changes, and no component
judges landslide risk. Node count, role mapping and disturbance parameters follow the frozen Task
v1.1 instance manifest; an earlier 16-node sketch is not used as a field fact.

Authority comes from two Chinese geological-hazard monitoring standards, DZ/T~0460--2023 and
DZ/T~0450--2023, retrieved and clause-audited in the repository~\cite{dzt0460,dzt0450}.
\S5.3.3 permits remote dynamic adjustment of sampling and upload frequency by warning level and
gives no numeric levels, so the periods used here are declared research choices; \S8.4.1 defines
four levels and \S8.4.2 permits upgrade, downgrade or termination after a conference; \S5.3.7
calls for dual-mode communication under weak coverage in high-risk extreme weather---a
conditional clause appearing once, not a blanket dual-channel mandate, and the communication-layer
standard does not impose it. The BeiDou open-service specification
(BDS-OS-PS-3.0)~\cite{bdsps3} reports $\geq$95\% success, sub-2~s average latency, an average
service frequency of one message per 30~s and a 14{,}000-bit single-message ceiling, noting that
realised frequency and maximum length are constrained by registration parameters: the rate is a
provisioned quantity. These clauses bound lawful action---re-compile frequency within granted
levels and request a conference change, never silently drop an obligation or declare risk.

The backup uses one 78-byte packet per 1200~s (20~B header, 58~B payload; a displacement sample is
6~B and rainfall 4~B, so roughly 9--10 samples fit). The elevated (yellow) task uses a 300~s
period with a delivery deadline two periods after window open, a $\approx$600~s delivery
window---shorter than the 1200~s backup slot spacing. The 200~B ceiling and $\geq$1~min period in
DZ/T~0450 (\S7.4.3.2, \S7.4.2.3, with a $\geq$5~s inter-message gap at \S7.4.3.5) belong to
different service-formulation sections and are not combined into one ``BeiDou operating point'';
the 1200~s/78~B point is a declared research choice.

\section{System Model}
\label{sec:model}

Time is slotted at $\Delta{=}60$~s and $t{=}0$ is 06:00 local, so $t{=}12$~h is sunset and the
solar cycle is explicit. Nodes $\mathcal N$, measurands $\mathcal M$ and a gateway $g$ compose the
field; a centre $c$ lies beyond the backhaul.

\subsection{Obligations and the delivery chain}
A monitoring obligation is $o=(n,m,[l_o,h_o),d_o)$ over node $n$, measurand $m$, sampling window
and absolute delivery deadline. A routine obligation at period $P$ has $h_o-l_o=P$ and
$d_o=h_o+P$, one period of grace. A sample fulfils $o$ only if it matches $(n,m)$, is taken
in-window and is received centrally by $d_o$, traversing
\begin{equation}
\resizebox{\columnwidth}{!}{$\displaystyle
\underbrace{\text{qualified sample}}_{\text{node}}
 \!\to\! \underbrace{\text{heard at }g}_{\text{LoRa access}}
 \!\to\! \underbrace{\text{returned by }d_o}_{\text{cellular / short message}}
 \!\to\! \underbrace{\text{central evidence}}_{\text{fulfilled}} .$}
\end{equation}
For piecewise missions the first segment uses a closed interval and later segments half-open
windows, so a boundary sample cannot double-cover adjacent windows after a level change.

\subsection{Links, outages, and who can act}
LoRa access (node$\to$gateway, with Class~A downlink windows) succeeds per opportunity with
probability $p_a{=}0.74$; the primary backhaul (gateway$\to$centre) is good with probability
$p_b{=}0.62$ outside a declared outage. A control outage is an interval $\mathcal O=[a,b)$ during
which the primary is unavailable to centre$\to$field traffic: centre commands and mission-table
segments are refused at admission and never transmitted, and the short-message backup cannot
carry downlink configuration. During $\mathcal O$ the only entities able to change field
communication behaviour are the \emph{node} (on its clock, own battery, cache, recent access
receipts and bounds carried by already-applied commands) and the \emph{gateway} (on samples it
has actually heard, its forwarding log and backup slots used). The centre retains authority to
authorise changes but no channel to install them; a revision authored in $\mathcal O$ reaches the
field at $b$.

\subsection{Persistent installed state}
\emph{Record state.} Each unacknowledged sample $s$ at node $n$ has creation time $t_s$, a
business deadline $d(s)=(\lfloor t_s/P\rfloor+2)P$ under the locally known cadence, and occupies
one of $K{=}240$ cache slots and up to $B{=}32$ positions in an access batch. Under the modelled
edge discipline it leaves the cache only on central acknowledgement (modelled synchronously
without airtime, equally favouring every queue rule); until then it is re-sent oldest-first.

\emph{Configuration state.} The node runs one profile $q(t)\in\{P_0,P_1\}$ (sparse
$P_0{=}600$~s, dense $P_1{=}300$~s) for both sampling and reporting, installed by an applied
command carrying a version. A command is \emph{applied} when the node receives and executes it;
a later report can confirm that effect to the observer. Queued or in-transit commands change
nothing. While applied,
$q{=}P_1$ draws additional sampling and uplink energy every tick until something terminates it.

\subsection{Energy}
A sample costs $e_s{=}4.7{\times}10^{-4}$~Wh and an uplink $e_u{=}2.33{\times}10^{-5}$~Wh;
battery capacity is $C{=}0.05$~Wh. Solar input is a 12~h half-sine with peak $\eta$ (nominally
$0.03$~Wh/h) plus per-node/hour stochastic cloud occlusion. A node whose state of charge reaches
zero \emph{permanently} dies in the main configuration. Holding dense over a 12~h night draws
$\approx0.071$~Wh $>C$.

\subsection{Local predicates and authorised fallback}
A record carries an expiry time $d_n(s)$ assigned from the cadence known at its source. Once
$t>d_n(s)$, retaining it has no value for obligations represented by that local contract.
Records remain eligible on the deadline tick. During an unreceived mission revision, the source
cannot substitute a future central deadline for $d_n(s)$; the independent scorer retains the
actual authorised obligations.

A configuration instead has a permitted fallback set and a local trigger, such as an installed
TTL, a clock rule or a battery condition. Stopping dense acquisition and deleting previously
collected records are separate actions: a sample can remain useful after the profile changes.
Consequently we do not assign configurations a universal release time derived from the last
return slot. Unknown future mission revisions cannot become node-local stopping evidence.

\subsection{Service loss and statement discipline}
The objective counts deadline-valid obligation delivery (service), permanent node deaths, and the
resource cost of state outliving its release time: wasted samples, access attempts, short-message
bytes and energy. Offline, each unfulfilled obligation is charged to its \emph{first} failing
stage under an optimistic model granting every later assumption, giving four mutually exclusive
segments: \senergy\ (no qualified sample), \saccess\ (sampled but not heard at the gateway by the
deadline), \stime\ (heard but no return opportunity exists) and \scap\ (an opportunity exists but
capacity or congestion cannot carry it). An online observer that cannot see a stage reports
\emph{unknown} and never guesses it. Every feasibility claim is \textbf{confirmed} (an independent
report shows the target config, or a matching sample is received centrally),
\textbf{unconfirmed} (a command is in flight or receipts are incomplete, which implies neither
that it took effect nor that the channel is down) or \textbf{unreachable} (the current path is
known down from explicit network-side state). A LoRa access receipt---``a node was heard''---is
evidence about the access link only and never by itself establishes backhaul reachability.

\section{Problem Formulation}
\label{sec:problem}

Let $\mathcal H_e(t)$ contain only the evidence actually available to entity $e$ by time $t$.
Node history includes clock, battery, cache and received commands; gateway history additionally
includes heard samples and its own forwarding log. A policy chooses record retention/release
and configuration hold/fallback through $a_e(t)=\pi_e(\mathcal H_e(t))$.
We compare the outcome vector
\begin{equation}
\mathcal J(\pi)=
\bigl(N_{\rm miss},N_{\rm brownout},E_{\rm node},A_{\rm up},A_{\rm down},B_{\rm backup}\bigr),
\label{eq:obj}
\end{equation}
where missed obligations, brownouts, node energy, uplink/downlink airtime and backup bytes are
reported separately. A scalar death penalty is a declared experimental preference in the
single-node reference, not an inferred requirement of disaster monitoring. Strict survival is
one possible risk regime, not a universal assumption imposed on every experiment.

\subsection{Execution contract}
The comparison obeys four conditions. \emph{Locality}: identical local histories cannot produce
different actions solely because of an unseen future. \emph{Authority}: configuration fallbacks
must lie within permissions granted before they are used. \emph{Retention accountability}: the
source may expire a record under its known lifetime; releasing transmission state does not erase
an obligation from the score. \emph{Outcome evidence}: a queued request, an applied configuration,
a gateway receipt and central delivery establish different facts.

These conditions leave a real tradeoff. A fallback may protect energy while reducing acquisition
for an outstanding obligation, and that loss remains in $N_{\rm miss}$. The system cannot promise
both full service and survival under every resource history merely by changing the controller.

\subsection{Where enforcement can act}
For a predicate $c$ and candidate layer $e$, the design condition is
\begin{equation}
\begin{split}
&\operatorname{evidence}(e,c)\ \land\ \operatorname{authority}(e,c)\\
&\qquad\land\ D_{\rm observe+decide+act}(e,c)<M_c(t),
\end{split}
\label{eq:placement}
\end{equation}
where $M_c(t)$ is the remaining intervention margin. Evidence and authority alone do not make a
central protection rule enforceable: its delay includes backhaul delivery and a Class~A window.
During a control outage that delay may exceed the energy margin. A source-local expiry needs
neither a new command nor a remote failure diagnosis. The condition identifies implementation
requirements; an empirical delay estimate is not a universal safety bound.

\subsection{Knowledge boundary}
Consider two histories with identical observations at the node, one with an unseen central
downgrade and one with the elevated requirement still active. Any causal node policy takes the
same action in both until new evidence arrives. It therefore cannot claim to track the unseen
mission revision exactly. A pre-authorised fallback resolves what the node is allowed to do;
it does not reveal which remote history occurred. This elementary indistinguishability argument
explains why an unknown outcome must remain distinct from confirmed completion or infeasibility.

An Agent, a deterministic rule and a solver can all consume this contract. Their common action
space and evidence boundaries define the communication problem independently of the planner.

\section{Mechanisms}
\label{sec:mech}
The prototype instantiates the execution contract at three locations. The source applies known
record lifetimes and authorised fallback rules; the gateway uses heard samples and finite return
opportunities; the centre compiles requirements and interprets received evidence. The experiments
below isolate these components. Their complete composition is not assumed to outperform an
ordinary stack containing the same mechanisms.

\subsection{Record termination: deadline-aligned expiry}
\label{sec:mech-record}
At each upload opportunity the node deletes cached records with $d_n(s)<t$ and sends survivors
oldest-first. We implement two formulations: \texttt{deadline-purge}, keyed to the obligation
cadence, and \texttt{generic-expiry}, an ordinary BPv7-style per-record lifetime assigned at
creation, $\text{expires}=t_s+((P-t_s\bmod P)+P)=d(s)$, referencing no obligation id, slot
geometry or certificate. A gateway variant locally acknowledges records it will not send. EDF and
obligation-greedy queues only reorder and never release; AoI \texttt{latest-only} releases
indiscriminately, keeping only the newest. Dropping stale or expired traffic is well established
in AoI packet management and DTN bundle expiry~\cite{costa2014pm,kam2016deadline,dtnarch,rfc9171},
and suppressing sensor traffic whose delay bound is exceeded is likewise prior
art~\cite{wada2009timeout}. We claim no new discard algorithm. The system-level content is the
end-to-end business-deadline basis---not packet age, hop TTL or an estimated single-packet
delay---zero-downlink source placement, and, decisively, cross-segment consistency, because the
acknowledgement coupling means a gateway-only release does not free the source cache.

\subsection{Configuration fallback with ordinary protection}
\label{sec:mech-lease}
The local clock guard clamps sampling and reporting to sparse at night and leaves daytime
commands untouched. The configuration comparison additionally evaluates fixed TTLs, a rolling
battery/expected-harvest check and an energy-derived stopping bound. TTL starts at the actual
over-the-air dense-application edge. All compared configurations use the same source expiry,
backup packing and permitted fallback profile.

Two information conditions are separated. Under announced validity, the node receives a known
validity end with the configuration. Under runtime authorisation, a downgrade issued during the
outage remains unknown in the field; local policies use their clock, measured battery and
declared energy model. Announced validity is an informative comparator, not a same-information
competitor to the latter policies. A separate quantised single-node stopping study compares
ordinary rules with an exact non-prescient reference under a public task table. Its scope and
findings are reported in \S\ref{sec:eval-lease}.

\subsection{Partial-observability projector}
\label{sec:mech-projector}
At compile time the centre holds the obligation table, the \emph{current} and never forecast
network-side outage state, and public parameters; at mission-table arrival the gateway holds only
what it has heard. A structural no-return certificate (\stime) is emitted from slot geometry alone
and is explicitly conditional on the primary being unavailable; geometry never infers an outage.
An arrival-time projector labels a missed obligation only when local evidence separates the
stages, and otherwise returns unknown. An unknown-aware centre projector emits \emph{unreachable}
solely from an explicit network-side field---which the current harness does not yet provide, so it
emits \emph{unconfirmed}---and replaces any blanket night kill with the computed energy ledger, so
genuine energy-infeasibility is distinguished from conservative advice.

\section{Prototype and Methodology}
\label{sec:proto}

All mechanisms run in one Python event-stepped simulator (60~s ticks) implementing LoRaWAN Class~A
downlink windows, versioned two-field configuration commands (sampling interval and reporting
period delivered as independent fields, with cross-generation mixed configs detected), the
hold-until-acknowledged edge cache, max-coverage backup packing, stochastic solar and cloud,
absorbing brownout, and the non-prescient piecewise mission-view gate. Disabling the backup
reproduces the v1.1 simulator bit-for-bit across five anchored invariants; the local guard and
fallback extensions are default-off.

Two controlled episodes are used. \emph{Episode~I}, the v1.1 manifest point, upgrades to yellow at
$t{=}6$~h with primary outage 4--20~h and the mission table reaching the gateway at 20~h; it drives
the frontier, attribution, record-expiry, placement and live-agent results. \emph{Episode~II} is a
deterministic configuration-termination scenario with identical parameters except schedule,
outage phasing and swept harvest: to place a daytime downgrade \emph{inside} an outage, phase~A
uses yellow at $t{=}2$~h with a conferred downgrade at $t{=}6$~h and outage 4--20~h, phase~B uses
yellow at $t{=}1$~h, downgrade at $t{=}8$~h and outage 6--22~h; $t{=}12$~h is sunset. Episode~II is
an explicitly constructed controlled comparison for the configuration sub-problem, not a new field
claim. Node deaths are reported as totals over seeds 0--2;
sub-point mean-service differences at $n{=}3$ are treated as not significant; seeds 6--9 are a
reused held-out confirmation set, not an untouched test set. Tables are generated from registered result files, and the artifact provides deterministic
checks and frozen reference results.

\section{Evaluation}
\label{sec:eval}

\subsection{Where the fixed-resource loss is}
\label{sec:eval-frontier}
Over Episode~I's 7560 routine obligations the strong maxcov/FIFO baseline delivers 3025 (0.400).
Symmetric per-obligation counterfactuals widen one physical resource at a time, never the policy,
and report set differences over identical obligation ids. Relaxing backhaul to unlimited gives
4740 (0.627), relaxing energy (0.50~Wh, always-dense) gives 4089 (0.541), and both gives 6023
(0.797); 1715 obligations are blocked by backhaul capacity alone, 1128 by battery alone, and 155
(2.0\%) by both (Table~\ref{tab:walls}). These are intervention results on implemented strategies,
not an exhaustive optimum; we make no zero-remainder or global-optimality claim. They locate substantial resource pressure at this operating point. The state-termination
experiments address an additional execution effect under fixed resources; their motivation does
not require proving that every scheduler has exhausted its headroom.

\begin{table}[t]
\centering\small
\caption{Counterfactual interventions over 7560 obligations (Episode~I). Set decomposition over the
tested configurations; not an upper bound over all legal policies.}
\label{tab:walls}
% 由 scripts/make_tables.py 生成，勿手改；数字来自结果文件。
\begin{tabular}{llcc}
\toprule
Run & Sampling/Energy & Backhaul & Delivered\\
\midrule
base & dayfeed, 0.05 Wh & 1200 s/78 B & 3025 (0.400)\\
$O_{bh}$ & dayfeed, 0.05 Wh & unlimited & 4740 (0.627)\\
$O_{en}$ & always-300, 0.50 Wh & 1200 s/78 B & 4089 (0.541)\\
$O_{bo}$ & always-300, 0.50 Wh & unlimited & 6023 (0.797)\\
\bottomrule
\end{tabular}

\end{table}

Charging each of the 4535 failures to its first failing segment, with ``heard'' required
\emph{by the deadline} rather than at any later time, gives \senergy\ 2002 (44.1\%), \scap\ 1115
(24.6\%), \saccess\ 830 (18.3\%) and \stime\ 588 (13.0\%). An earlier heard-ever accounting had
reported \scap\ 1945 and \saccess\ 0: all 830 moved obligations were in fact sampled but first
heard at the gateway after their deadline, late access rather than backhaul capacity. The
delivered count (3025) is identical under both accountings; the correction changes \emph{where the
loss is charged}, and therefore which mechanism could address it.

\subsection{Record expiry: standard primitive, decisive placement}
\label{sec:eval-record}
During the Episode~I outage 75 backup packets carry 750 records; under FIFO 219 arrive centrally
already expired against 261 on-time (45.6\%; 1270 of 2782 payload bytes). Gateway-only suppression
removes those sends but withholds the acknowledgement, and on-time delivery falls $202\to194$: cost
moves to the access segment. Source-local deadline expiry reverses the sign, raising outage on-time
delivery $\cThreeSeedZeroFifo\to\cThreeSeedZeroPurge$ Across ten seeds
(Table~\ref{tab:purgeexpiry}) expiry raises mean service $\cThreeSvcFifo\to\cThreeSvcPurge$, a paired $+\cThreePairedPoints$ points
(95\% CI $+\cThreePairedLo$ to $+\cThreePairedHi$, positive on all ten), total outage on-time delivery $\cThreeTenSeedFifo\to\cThreeTenSeedPurge$
(\cThreeRatio$\times$), expired backup records $\cThreeExpiredFifo\to\cThreeExpiredPurge$, and deaths $\cThreeDeathsFifo\to\cThreeDeathsPurge$; it beats AoI latest-only by
$+\cThreeLatestDelta$ points, because latest-only buys outage freshness by abandoning
post-recovery coverage. Crucially, \texttt{deadline-purge} and \texttt{generic-expiry} are
bit-identical on every seed---service, outage counts, deaths and delivered set---so the record
mechanism \emph{is} standard per-record expiration, and its measured value is the business-deadline
basis plus source/gateway cross-segment consistency, not a new discard algorithm.

\begin{table}[t]
\centering\small
\caption{Ten-seed Episode~I sweep with the strong maxcov chooser fixed; only the edge disposition
changes. Gateway-only suppression (not shown; same FIFO cache) gives 194 outage on-time at seed~0,
against \cThreeSeedZeroFifo for FIFO and \cThreeSeedZeroPurge for source expiry---the cross-segment cost. $^{\dagger}$standard BPv7-style per-record lifetime; $^{\ddagger}$obligation-keyed purge, bit-identical to $^{\dagger}$.}
\label{tab:purgeexpiry}
\setlength{\tabcolsep}{3.5pt}\footnotesize
% 由 scripts/make_tables.py 生成，勿手改；数字来自结果文件。
\begin{tabular}{lcccc}
\toprule
Edge disposition & Full svc & Outage OT & Expired & Deaths\\
\midrule
FIFO (baseline) & 0.373 & 1691 & 2538 & 12\\
latest-only (AoI) & 0.358 & 4250 & 0 & 0\\
\texttt{generic-expiry}$^{\dagger}$ & 0.415 & 4265 & 1 & 0\\
\texttt{deadline-purge}$^{\ddagger}$ & 0.415 & 4265 & 1 & 0\\
\bottomrule
\end{tabular}

\end{table}

\subsection{Configuration protection and its algorithmic boundary}
\label{sec:eval-lease}
Episode~II holds record expiry and maxcov packing fixed while comparing local fallback rules.
The clock-only rule reacts at sunset; fixed-duration and energy rules can act earlier. The
announced-validity row has advance mission-end information and is shown separately from policies
that cannot observe a downgrade during the outage.

\begin{table*}[t]
\centering\small
\caption{Configuration fallback across two outage phases (\cFiveSeeds\ seeds,
\texttt{generic-expiry} and maxcov fixed). TTLs start at dense application. The announced-validity
row has Task~1 information; the remaining adaptive/TTL policies use Task~2 information.
The three reported energy variants have identical cell outcomes, so one row represents them.
Zero observed deaths is not a universal survival guarantee.}
\label{tab:lease}
% 由 scripts/make_tables.py 生成，勿手改；数字来自结果文件。
\begin{tabular}{lcccc}
\toprule
Local revert bound & Deaths (3 seeds) & Mean svc & Mean final SoC & Yellow delivered / total\\
\midrule
\multicolumn{5}{l}{\emph{Phase A, overcast $\eta{=}0.012$, seeds 0--2 (upgrade 2~h, downgrade 6~h, outage 4--20~h)}}\\
never dense (blue only) & 0 & 0.5752 & 0.00610 & 547/2016\\
clock guard (sunset 12~h) & 20 & 0.5961 & 0.00329 & 849/2016\\
announced validity (6~h, Task~1) & 0 & 0.5970 & 0.00610 & 849/2016\\
fixed TTL 4~h (revert 6~h) & 0 & 0.5970 & 0.00610 & 849/2016\\
fixed TTL 8~h (revert 10~h) & \textbf{0} & \textbf{0.5971} & 0.00552 & \textbf{849/2016}\\
energy gate, open-loop bound & \textbf{0} & \textbf{0.5771} & 0.00504 & \textbf{559/2016}\\
\midrule
\multicolumn{5}{l}{\emph{Phase B, overcast $\eta{=}0.012$, seeds 0--2 (upgrade 1~h, downgrade 8~h, outage 6--22~h)}}\\
never dense (blue only) & 0 & 0.5399 & 0.00609 & 1168/3528\\
clock guard (sunset 12~h) & 31 & 0.5982 & 0.00191 & 2020/3528\\
announced validity (8~h, Task~1) & 0 & 0.5994 & 0.00609 & 2020/3528\\
fixed TTL 4~h (revert 5~h) & 0 & 0.5928 & 0.00606 & 1916/3528\\
fixed TTL 8~h (revert 9~h) & \textbf{0} & \textbf{0.5996} & 0.00606 & \textbf{2020/3528}\\
energy gate, open-loop bound & \textbf{0} & \textbf{0.5431} & 0.00416 & \textbf{1194/3528}\\
\bottomrule
\end{tabular}

\end{table*}

At $\eta=\cFivePeak$, the clock guard leaves \cFiveNightfloorDeadA\ and
\cFiveNightfloorDeadB\ deaths in phases A and B. Fixed TTL 8~h yields zero deaths and yellow
delivery counts $\cFiveTtlEightYellowA/\cFiveYellowNA$ and
$\cFiveTtlEightYellowB/\cFiveYellowNB$, whereas TTL 4~h loses
\cFiveTtlFourLostB\ deliveries in phase B. The evaluated energy-derived gate also yields zero
deaths but loses \cFiveEnergyLostA\ and \cFiveEnergyLostB\ yellow deliveries relative to TTL
8~h. Its predictive load approximation is conservative relative to the measured radio load;
these results characterise that implementation, not all predictive energy control.
The matrix establishes the utility of local fallback while providing no independent increment
for the proposed energy-derived lease.

\textbf{Same-information stopping reference.} To separate an implementation limitation from
optimisation headroom, a quantised single-node model evaluates continued dense operation versus
an absorbing return to sparse over the full physical horizon. Its public task table determines
reward; cloud realisations remain unknown. The model is calibrated to the current simulator's
sampling and radio load, and ordinary TTL, battery-threshold and rolling rules receive the same
information. Across \cNineCells\ phase/harvest/variability cells and \cNineWeights\ declared death
penalties, ordinary combinations match the exact non-prescient reference in service in
\cNineServiceMatches/\cNineComparisons\ comparisons and in the priced objective up to numerical
precision. The strict all-futures survival endpoint is infeasible at the declared tight-energy
point, including for all-sparse operation. This is a risk-model boundary, not a claim that the
monitoring mission or every physical implementation is infeasible.

The negative result closes the tested single-node stopping family. It neither supplies a new
configuration algorithm nor establishes optimality for shared network scheduling or unseen task
tables. Standard fallback remains a component of the execution system.

\subsection{What the field can attribute when the table arrives}
\label{sec:eval-attr}
In Episode~I the delayed mission table reaches the gateway at $t_a{=}20$~h, by which point 2136
yellow obligations are past deadline. Offline time-aware attribution labels all 2136 correctly
(\stime\ 588, \scap\ 327, \saccess\ 650, \senergy\ 571). Using only evidence the gateway lawfully
holds by $t_a$ labels 1138/2136 (53.3\%)---all geometric \stime\ (588), all in-time-not-returned
\scap\ (327) and 223 provably late \saccess---at 100\% precision on the labelled subset; 998
(46.7\%) remain unknown, whose offline truth is 571 never-sampled and 427 heard only after $t_a$
(Table~\ref{tab:attribution}). ``Never sampled'' and ``sampled but not yet heard'' are inseparable
without evidence the gateway does not have, and must stay unknown, which is the original G5
constraint. Compile-time geometry, needing no sample, emits 602 conditional no-return certificates
of which 588 are true \stime, the 14 recovery-edge cases being excluded rather than certified; an
optimistic-energy capacity certificate is impure (552 correct \scap, 571 actually \senergy, 193
false positives), so capacity is left to the gateway. The design consequence is that expiry under a known local deadline can proceed while the
remote cause remains unresolved; configuration fallback separately requires prior authority.

\begin{table}[t]
\centering\small
\caption{Arrival-time attribution of the 2136 upgraded obligations already missed when their table
reaches the gateway (Episode~I). The online projector is exact but partial; 427 matching samples
are first heard after $t_a$ and 650 after their deadline.}
\label{tab:attribution}
\setlength{\tabcolsep}{3.5pt}\footnotesize
% 由 scripts/make_tables.py 生成，勿手改；数字来自结果文件。
\begin{tabular}{lcc}
\toprule
Segment & Offline truth & Online at $t_a$\\
\midrule
\stime (no return slot; geometry) & 588 & 588\\
\scap (heard by deadline, not returned) & 327 & 327\\
\saccess (heard after deadline) & 650 & 223\\
\senergy (no qualified sample) & 571 & --- (unknown)\\
\unk (not separable online) & --- & 998\\
\midrule
Total & 2136 & 2136\\
Labelled coverage & 100\% & 53.3\%\\
Precision on labelled subset & 1.000 & 1.000\\
\bottomrule
\end{tabular}

\end{table}

\subsection{Enforcement placement determines survival}
\label{sec:eval-place}
A centre-side feasibility \emph{envelope}---a pure filter over any planner using only confirmed
configs, hearing counts, the authorised requirement and the clock---removes the command storm but
cannot save the tightest seed, because a dusk downgrade must cross the backhaul and a Class~A
window it misses. The same ``never stay dense overnight'' rule is a late, refusable suggestion at
the centre and an enforceable local clamp at the node (Table~\ref{tab:r39}). Naive centre
compliance kills 13 of 14 nodes on all three seeds (39 total) at service 0.162; adding the
node-local guard removes all deaths at unchanged service; with the guard present the envelope
additionally cuts refused admission attempts to 222. A zero-command pre-provisioned rhythm reaches
0.360 service, matching every centre-compiled arm for these static missions---a scoped negative
result on centre sampling control for pre-provisioned tasks. The ordinary local controller is
therefore part of the baseline for any further runtime or Agent comparison.

\begin{table*}[t]
\centering\small
\caption{Deterministic Episode~I placement study, no LLM. The local guard removes all permanent
deaths in the tested seeds at unchanged service; the centre envelope alone cannot. Refused commands
are rejected at admission and never transmitted, which is command hygiene rather than billing
saving. Point-specific empirical results, not a survival invariant or optimal bound.}
\label{tab:r39}
% 由 scripts/make_tables.py 生成，勿手改；数字来自结果文件。
\begin{tabular}{lcccc}
\toprule
Configuration & Mean svc & Deaths & Cmds sent & Refused\\
\midrule
comply (naive, unprotected) & 0.162 & 39 & 160 & 2488\\
dayfeed-c (centre day/night rule) & 0.370 & 12 & 308 & 1256\\
envelope(comply), centre only & 0.359 & 12 & 308 & 298\\
comply + local guard & 0.366 & 0 & 1424 & 3364\\
dayfeed-c + local guard & 0.368 & 0 & 258 & 1156\\
\textbf{envelope(comply) + local guard} & \textbf{0.359} & \textbf{0} & \textbf{282} & \textbf{222}\\
\midrule
\textbf{pure-local (pre-provisioned, 0 cmds)} & \textbf{0.360} & \textbf{0} & \textbf{0} & \textbf{0}\\
\bottomrule
\end{tabular}

\end{table*}

\subsection{Formative real-agent study}
\label{sec:eval-agent}
We ran deepseek-flash live over the Episode~I outage: three seeds each of an unaided agent (A0), the
agent with an ordinary pending-command checklist (A0s) and a certificate agent (A1), plus two
healthy-link controls, with a real model call at each 30-minute decision, 99 per trace and 1089
decisions in total; all arms share harness, legal observation, tools, history and action space.
The study motivates and stress-tests the interface and is not used to establish the deterministic
mechanisms.

\emph{Service and survival, noisy at $n{=}3$.} A1 attains mean service 0.358, level with the 0.360
deterministic safe floor, with 4 deaths; A0 0.296 with 35 deaths and A0s 0.222 with 39; the
healthy-link controls give A0 0.275/13 and A1 0.511/0. Per-seed service overlaps, so we draw no
service significance. The live agent loop was not wired to the local guard, so the four A1 seed-1
deaths are not cleared by the deterministic result of \S\ref{sec:eval-place}.

\emph{Call accounting.} The 1089 decisions are not 1089 successful calls: 52 ended in a
parse-failure hold from \texttt{max\_tokens} cutting the JSON mid-string. Per-arm parse-failure rates
are A0 7.1\% (28/396), A0s 6.4\% (19/297) and A1 1.3\% (5/396), so arm differences mix
output-format reliability with any semantic effect.

\emph{The interface fault, symmetrically scored.} The observation's \texttt{link} field carries
only LoRa access receipts (\texttt{n\_heard}) and no backhaul state. In the forced outage window
A0/A0s assert ``link up, commands can apply'' on 36/32 decisions, and A1 on 1 under a
token-aware statement check---over-optimistic, since the primary
is down throughout, and accompanying the ineffective dense orders and deaths. In the opposite
direction the v4 certificate asserted ``control plane not delivering'' whenever fewer than all
nodes were heard, firing at five healthy-link times when the primary was in fact usable, twice
after 13 of 14 nodes had already confirmed dense. Statements are checked in both outage and
healthy conditions, independently of whether the decision issued a command.
Offline replay of an unknown-aware v5 projector on identical snapshots issues zero unreachable
assertions without a network-side field, turning all five healthy false alarms into unconfirmed,
and splits 528 blanket night kills into 189 genuinely energy-infeasible and 339 advisory, every
near-dawn 04:30--06:00 case being advisory. The finding is that a missing backhaul-state
signal causes bidirectional errors; isolating any independent increment of the projector requires a
matched-capability arm receiving the same constants and heuristics without the projector,
end-to-end v5 runs with an explicit backhaul-state field, and cross-model replication, all future
work.

\section{Related Work}
\label{sec:related}

\textbf{LLMs in emergency and network control.} TopoLLM applies adaptive tool learning to
post-disaster topology planning~\cite{topollm}, and later systems let LLMs schedule UAV paths and
bandwidth, shape RL rewards or orchestrate recovery tools, including intent objects, budget and
time constraints, candidate-plan checks, local repair and documented network-management
APIs~\cite{icgrestore,ossgpt}. We claim no novelty for natural language to constraints, nor for
tool orchestration. Our comparison concerns a different execution setting: persistent field state, Class~A
actuation and a separately interrupted return path. We use those works to position the controller
interface and do not claim to have reproduced or refuted their original systems.

\textbf{Scheduling, freshness, energy.} The closest state-model precedent is an LLM in-context
data-collection scheduler with bounded per-sensor queues, exogenous arrivals, finite service
opportunities and feedback~\cite{icldc}; AoI and energy-harvesting control, including
instantaneous-battery threshold policies and partially observed harvesting, address resource
pieces~\cite{bacinoglu1905,zakeri2311,costa2014pm}; overbooking and contact-aware forwarding for
space links~\cite{ccsdssabr} and network-calculus bounds~\cite{leboudec01} are complementary
deterministic characterisations. Joint sampling, transmission and energy optimisation is
established and not claimed here; our distinction is persistent \emph{remote} configuration with
Class~A apply latency, a separated access versus primary/backup return path, and termination under a
control outage. Consistent network update research has long studied per-update validity,
dependencies and safe transitions~\cite{sdnconsistency}; we rely on that lesson and add the
disconnected, energy-constrained edge case in which an installed update's \emph{own data} become
undeliverable before a counter-update can arrive.

\textbf{Packet discard, DTN expiry, contacts.} AoI management pre-empts or drops old
updates~\cite{kam2016deadline,krishnan2019multihop}; deadline-aware buffer discard has been coupled
with adaptive ARQ for video~\cite{razavivideo}; DTN bundle protocols attach creation-time lifetimes
and use TTL or encounter drop~\cite{dtnarch,iranmaneshdtnqueue,rfc9171}, and custody
transfer~\cite{rfc5050} moves retention responsibility as a forward hand-off, since migrated out of
the BPv7 core. Rural LoRa--LEO systems schedule around estimated satellite contacts and exploit
spatial gateway diversity~\cite{sateriot}. Our record disposition \emph{is} BPv7-style expiry and we
demonstrate equivalence on the tested seeds; ordinary TTL and battery protection also cover our evaluated configuration-stopping family.
The measured system contribution is the placement and end-to-end effect of these primitives in
the declared monitoring model. Custody transfer is a substantive alternative to central-ACK
retention and must remain a comparator when feedback and storage responsibilities are changed.

\section{Discussion, Limitations, and Threats to Validity}
\label{sec:disc}

\textbf{What the results establish.} Source expiry, local protection and partial attribution have
separate evidence. Their common lesson is to place decisions where the required evidence and
ability to act survive disconnection. The strongest positive result concerns source retention;
configuration stopping is a boundary study showing where ordinary mechanisms suffice. A better
scheduler would not invalidate the observed retention or unreachable-downgrade failures.

\textbf{Composition and generalisation.} A complete execution runtime must be compared against
an ordinary stack with the same expiry, local protection, configuration shadow, observations and
planning tools. Component effects cannot be summed to infer an independent runtime gain. The
present results cover a declared deployment and controlled episodes; outage phase, coupled
access/backhaul failures and energy conditions delimit transfer. Firmware support for local
fallback and expiry metadata, their signalling/storage cost, and hardware execution have not
been established by the simulation results.

\textbf{Feedback and retention.} Central acknowledgement is synchronous and airtime-free in the
model, equally favouring the compared queue rules. With an uplink-only backup, actual central
delivery need not immediately become source knowledge. A deployed design must distinguish a
local access receipt, durable gateway custody and central delivery, and account for any release
notification sent through Class~A windows. Standard custody transfer, cumulative acknowledgements
and source expiry are therefore essential ordinary comparators. Joint allocation of finite
feedback and configuration capacity remains outside the evaluated model.

\textbf{Mission and energy scope.} During Episode~I, the source assigns expiry using the locally
known pre-change cadence; central scoring still uses the actual mission. Multiple overlapping
obligations may require retaining a record until its last valid use. Harvesting is synthetic and
brownout is absorbing in the main configuration. The exact stopping reference has a public task
table and a declared quantised energy model; its optimality statement does not transfer to an
unknown mission or a complete communication network.

\textbf{Agent evidence and remaining interface question.} The live study uses one model and
has both output-format and expert-information confounds. Its supported result is the interface
failure and the projector's offline correction. A matched-capability end-to-end comparison must
share local protection, evidence, deterministic tools and parsing budgets. The relevant outcome is
whether an action-aware interface reduces harmful or ineffective execution while retaining
beneficial changes, along with independently checked completion statements. Cross-model
replication follows an isolated effect; more calls alone cannot establish it.

\textbf{Provenance.} Standard clauses were checked against the repository's source copies;
one copy chain uses a document-sharing host because the official endpoint was unavailable.
The instance parameters distinguish published capabilities from research choices. Cited systems'
implementations were not executed, and no reproduction of their reported performance is claimed.

\section{Conclusion}
Intermittent backhaul can remove control while installed communication state continues consuming
resources. In the studied pre-disaster monitoring network, standard source expiry improves
on-time service by \cThreePairedPoints\ percentage points and outage delivery by
\cThreeRatio$\times$; local protection acts where a later central downgrade may be unreachable.
Ordinary rules cover the evaluated configuration-stopping problem. The resulting design principle
is to align retention and fallback with locally available evidence, authority and enforceable
actuation, while keeping end-to-end completion distinct from intermediate receipts. An Agent can
use this execution interface, but its independent benefit requires a matched ordinary-stack
comparison. The present evidence supports the communication mechanisms and their boundaries,
without making planner intelligence a prerequisite for correct local action.

\bibliographystyle{IEEEtran}
\bibliography{../refs}

\end{document}
